\subsection{Localisation des révisions de l’assemblage humain}\label{sec:paysage-chromosomique}
Pour situer les révisions entre les versions v1.0 et v1.1 de CHM13, on relève l’intervalle de référence \([a_i,b_i)\) de chaque entrée du VCF du consortium~\cite{chmreleases,mccartney}. Sur un chromosome de longueur \(L_i\), la plus petite marge terminale qui contient l’intervalle entier a pour largeur \(r_i=\min(b_i,L_i-a_i)\). La figure~\ref{fig:paysage-chromosomique} cumule les entrées en fonction de cette largeur,
\[
 F_w(h)=\frac{\sum_i w_i\,\mathbf{1}_{r_i\le h}}{\sum_i w_i}.
\]
Le poids \(w_i=1\) compte chaque entrée une fois. Le poids \(w_i=d(\mathrm{REF}_i,\mathrm{ALT}_i)\) mesure le coût du remplacement de son allèle. Le total de 2\,663 éditions couvre les allèles consignés dans le VCF. Les quatre substitutions résiduelles du rejeu interviennent séparément dans la borne de la section~\ref{sec:panelhumain}.

\chart{biology/paysage_chromosomique.pdf}{Le panneau a) donne la part des 310 entrées et de leur coût couverte par des marges de 0 à 1\,000~nt. Le panneau b) répartit ce coût entre les deux marges de 500~nt et les autres entrées de chaque chromosome. Les cases blanches ont un coût nul.}{fig:paysage-chromosomique}{python/studies/chromosome_landscape.py}
\FloatBarrier
\Needspace{4\baselineskip}
La marche à 492~nt correspond à une délétion de 491 bases au début du chromosome~3. Cet allèle représente à lui seul 18,4\,\% du coût du VCF sur le panel. À l’autre bout de la distribution, une substitution sur le chromosome~9 exige une marge de plus de 16~Mnt. Le découpage des chromosomes doit donc aussi tenir compte de révisions intérieures.

\synthese{Sur les 18 chromosomes, les 500 bases de chaque extrémité couvrent 18\,000 des 2\,643\,920\,725 bases de référence. Elles contiennent 263 des 310 entrées, soit 84,8\,\% des entrées et 98,2\,\% de leur coût.}
